% Suggested insertion in the proof of A:hinich-explicit.
% Independent audit by /root/algebra_c, 19 November2026.
\paragraph{An independent explicit cross-ratio verification.}
For these particular coordinates one can compute the entire plumbing family directly, without using the cotangent pairing. The first capped component has a global coordinate $\zeta_1$: on the original bounded component it is the inverse of the injective map $w=\phi(\zeta_1)$, and it is the exterior coordinate on its attached cap. Its three marks are $A,B,\infty$. The second capped component has the original global $w$-coordinate: on its attached interior cap use $w=\phi(\zeta_2)$. Its three marks are $c,\infty,R^2$. The two node coordinates are therefore exactly
\[
 x=(1-\varepsilon)/\zeta_1,
 \qquad y=\zeta_2/(1+\varepsilon).
\]

Fix any $t$ with $0<|t|<1$, and put $\alpha=t/t_0$. The first retained plumbing piece is
\[
 |\zeta_1|<(1-\varepsilon)/|t|,
\]
and the second retained piece is the complement in the $w$-sphere of
\[
 \phi\bigl(\{|\zeta_2|\le(1+\varepsilon)|t|\}\bigr).
\]
On the first piece define $w_t=\phi(\alpha\zeta_1)$; on the second use its existing $w$-coordinate. The plumbing equation $xy=t$ is equivalent to $\zeta_2=\alpha\zeta_1$, so these functions agree on the overlap. Because
\[
 |\alpha|\,(1-\varepsilon)/|t|
       =(1-\varepsilon)/t_0=1+\varepsilon<R,
\]
and $\phi$ is injective on $|\zeta|<1+\varepsilon$, the first piece maps biholomorphically to $\phi(\{|\zeta|<1+\varepsilon\})$. Its overlap with the second piece is exactly the image of
\[
 (1+\varepsilon)|t|<|\zeta_2|<1+\varepsilon.
\]
The two images cover the $w$-sphere. Thus $w_t$ is a global conformal coordinate on the plumbed sphere for every $0<|t|<1$, not merely near $t_0$.

The four labelled marks in this coordinate are
\[
 a(t)=\phi((t/t_0)A),\quad
 b(t)=\phi((t/t_0)B),\quad c=\phi(C),\quad\infty.
\]
They stay distinct throughout this punctured unit disc: the first two lie in the injective $\phi$-disc, while $c,\infty$ are outside it. A cross-ratio coordinate on labelled $M_{0,4}$ is therefore the explicit rational function
\[
 \Lambda(t)=\frac{\phi((t/t_0)A)-c}{\phi((t/t_0)B)-c}.
\]
Substituting the specified rational constants and differentiating gives
\[
 \Lambda'(t)=\frac{201579840}{323761}\,
 \frac{(101t-99)(101t+12518)}
 {(101t-1320)^2(101t+924)^2}.
\]
In particular
\begin{align*}
 \Lambda(t_0)&=\frac{323680}{323761},
 &\Lambda'(t_0)&=0,\\
 \Lambda'(0)&=-\frac{909}{1752520}\ne0,
 &\Lambda''(t_0)&=\frac{22850240}{44933817907}\ne0.
\end{align*}
Thus the nonzero parameter $t_0$ is a simple critical point, while the family is a local smoothing coordinate at the nodal parameter $0$. This verifies the geometric conclusion independently of the Kodaira--Spencer computation. The exact certificate \texttt{code/C\_plumbing\_crossratio.py}, with output \texttt{logs/C\_plumbing\_crossratio.json}, also checks that all zeros and poles of $\Lambda$ lie outside the closed unit disc, and that $\Lambda=1$ occurs there only at $t=0$. Its eleven assertions use exact rational functions.
